To compute $s_q(\text{GL}_2(p))$—the number of conjugacy classes of subgroups of order $q$ in $G = \text{GL}_2(p)$—for an odd prime $q \neq p$, we analyze the structure and actions of these subgroups. 

For an odd prime $q$, any subgroup of order $q$ is cyclic (isomorphic to $C_q$). The formula given in Proposition 3.2(a) is:
$$s_q(\text{GL}_2(p)) = \frac{q+3}{2} \Delta^q_{p-1} + \Delta^q_{p+1}$$
where $\Delta^a_b = 1$ if $a$ divides $b$, and $0$ otherwise. 

Since $q$ is an odd prime, it cannot divide both $p-1$ and $p+1$ simultaneously (otherwise it would divide their difference, $2$, which is impossible for an odd prime). Thus, the divisibility conditions are mutually exclusive, leading to two distinct cases.

---

### Case 1: $q \mid (p+1)$ (meaning $q \nmid (p-1)$)
In this case, $\Delta^q_{p-1} = 0$ and $\Delta^q_{p+1} = 1$. The formula predicts $s_q(\text{GL}_2(p)) = 1$ class.

* **Algebraic Explanation:**
  Since $q \nmid (p-1)$, the field $\mathbb{F}_p$ does not contain any primitive $q$-th roots of unity. Consequently, there are no reducible (diagonalizable) subgroups of order $q$ because any such subgroup would require eigenvalues in $\mathbb{F}_p$.
  
  However, since $q \mid (p+1)$, it must divide $p^2 - 1 = (p-1)(p+1)$. Thus, the field extension $\mathbb{F}_{p^2}$ contains primitive $q$-th roots of unity. In $\text{GL}_2(p)$, the multiplicative group $\mathbb{F}_{p^2}^*$ embeds as a **Singer cycle** $S \cong C_{p^2-1}$. Since $S$ is cyclic, it contains a unique subgroup of order $q$. Because its eigenvalues lie in $\mathbb{F}_{p^2} \setminus \mathbb{F}_p$, this subgroup $H$ is **irreducible**.
  
  Any irreducible subgroup of order $q$ is conjugate to a subgroup of the Singer normalizer $N \cong C_2 \ltimes S$. Since $S$ is normal in $N$ and has index $2$, and $q$ is an odd prime, the unique subgroup of order $q$ in $S$ is the only subgroup of order $q$ in $N$. Thus, there is exactly **$1$ conjugacy class** of subgroups of order $q$.

---

### Case 2: $q \mid (p-1)$ (meaning $q \nmid (p+1)$)
In this case, $\Delta^q_{p-1} = 1$ and $\Delta^q_{p+1} = 0$. The formula predicts $s_q(\text{GL}_2(p)) = \frac{q+3}{2}$ classes.

* **Algebraic Explanation:**
  Since $q \nmid (p^2-1)$, there are no irreducible subgroups of order $q$. Every subgroup of order $q$ is **reducible** and can be simultaneously diagonalized over $\mathbb{F}_p$. Thus, up to conjugacy, any subgroup $U$ of order $q$ is a subgroup of the diagonal group $D \cong C_{p-1} \times C_{p-1}$.

  Let $a \in \mathbb{F}_p^*$ be an element of order $q$. Any non-trivial subgroup of order $q$ inside $D$ is cyclic and can be generated by a diagonal matrix of the form $\text{diag}(a^i, a^j)$ where $(i, j) \not\equiv (0, 0) \pmod q$. 

  We can classify these subgroups by using a parameter $\ell \in \mathbb{F}_q \cup \{\infty\}$ representing the "slope" or ratio of the powers of the diagonal entries:
  1. If the first diagonal entry is $1$ (i.e., $i \equiv 0$), we get the subgroup:
     $$U_{\infty} = \langle \text{diag}(1, a) \rangle$$
  2. If the first diagonal entry is not $1$ (i.e., $i \not\equiv 0$), we can rescale the generator by raising it to the power $i^{-1} \pmod q$. This yields a unique generator of the form $\text{diag}(a, a^{\ell})$ for some $\ell \in \mathbb{F}_q$:
     $$U_{\ell} = \langle \text{diag}(a, a^{\ell}) \rangle$$

  To find the number of conjugacy classes, we look at how these subgroups partition under conjugation. The normalizer of $D$ in $\text{GL}_2(p)$ is the monomial group $I = C_2 \ltimes D$, where the $C_2$ generator swaps the diagonal entries:
  $$\text{diag}(x, y) \mapsto \text{diag}(y, x)$$

  Swapping the entries of $U_{\ell} = \langle \text{diag}(a, a^{\ell}) \rangle$ produces $\langle \text{diag}(a^{\ell}, a) \rangle$. Raising this generator to the power $\ell^{-1} \pmod q$ normalizes the first entry to $a$, giving $\langle \text{diag}(a, a^{\ell^{-1}}) \rangle = U_{\ell^{-1}}$.
  
  Therefore, the conjugation action acts on the parameter space $\mathbb{F}_q \cup \{\infty\}$ via the map $\ell \mapsto \ell^{-1}$. We count the orbits of this action:
  * **Orbit 1:** $\{0, \infty\}$. The subgroups $U_0 = \langle \text{diag}(a, 1) \rangle$ and $U_{\infty} = \langle \text{diag}(1, a) \rangle$ are conjugate. (**1 class**)
  * **Orbit 2:** $\{1\}$. Since $1 \equiv 1^{-1} \pmod q$, the subgroup $U_1 = \langle \text{diag}(a, a) \rangle$ (the scalar matrices of order $q$) forms its own orbit. (**1 class**)
  * **Orbit 3:** $\{-1\}$ (or $q-1$). Since $-1 \equiv (-1)^{-1} \pmod q$, the subgroup $U_{-1} = \langle \text{diag}(a, a^{-1}) \rangle$ forms its own orbit. (**1 class**)
  * **Remaining Orbits:** The remaining $q-3$ elements of $\mathbb{F}_q^* \setminus \{1, -1\}$ do not equal their own inverses. They pair up into orbits of size 2 of the form $\{\ell, \ell^{-1}\}$. This yields:
    $$\frac{q-3}{2} \text{ classes}$$

  Summing these classes gives:
  $$\text{Total Classes} = 1 + 1 + 1 + \frac{q-3}{2} = 3 + \frac{q-3}{2} = \frac{q+3}{2}$$
  This matches the coefficient in the formula.